\documentclass[10pt]{article}
\usepackage[a4paper, margin=1in]{geometry}
\usepackage{amsmath}
\usepackage{amssymb}
\usepackage[hangul]{kotex}
\usepackage{bm}
\usepackage{xcolor}
\usepackage{algorithm}
\usepackage{algpseudocode}

\title{\textbf{SWOMP 알고리즘 분석 및 OMP와의 비교}}
\author{Gemini}
\date{\today}

\begin{document}
\maketitle

\section{서론}
하이브리드 빔포밍(Hybrid Beamforming)은 대규모 MIMO(Massive MIMO) 시스템의 하드웨어 복잡도를 낮추기 위한 핵심 기술입니다. 이 구조에서 최적의 완전 디지털 빔포머($F_{opt}$)를 아날로그 빔포머($F_{RF}$)와 디지털 빔포머($F_{BB}$)로 근사하는 것은 매우 중요한 문제입니다. OMP(Orthogonal Matching Pursuit) 계열 알고리즘은 이 문제를 해결하기 위한 효과적인 접근법을 제공합니다. 본 문서는 표준 OMP의 의사 코드(Pseudo-algorithm)와 그 한계를 분석하고, 이를 개선한 SWOMP(Simultaneous Weighted OMP)의 의사 코드를 정리하며, 두 알고리즘 간의 핵심 차이점을 명확히 부각시키는 것을 목표로 합니다.

\section{표준 OMP (Orthogonal Matching Pursuit)}
표준 OMP는 모든 데이터 스트림(고유 모드)을 동등한 중요도로 간주하고, 전체적인 근사 오차를 최소화하는 데 집중합니다.

\subsection{목표 함수}
Frobenius Norm으로 정의된 근사 오차를 최소화합니다.
\begin{equation}
    \min_{F_{RF}, F_{BB}} ||F_{opt} - F_{RF}F_{BB}||_F^2
\end{equation}

\subsection{의사 코드}
\begin{algorithm}[H]
    \caption{Standard OMP for Hybrid Beamforming}
    \begin{algorithmic}[1]
        \State \textbf{입력:} $F_{opt} \in \mathbb{C}^{N_t \times N_s}$, 코드북 $\mathcal{A} \in \mathbb{C}^{N_t \times N_c}$, RF 체인 수 $N_{RF}$
        \State \textbf{초기화:} 잔여 행렬 $R \gets F_{opt}$, 아날로그 빔포머 $F_{RF} \gets [\ ]$
        \For{$k = 1$ to $N_{RF}$}
            \State \textbf{매칭(Matching):} $R$과의 상관관계가 가장 큰 빔을 코드북에서 선택
            \begin{equation*}
                \mathbf{a}_k \gets \arg\max_{\mathbf{a} \in \mathcal{A}} ||\mathbf{a}^H R||_F^2
            \end{equation*}
            \State \textbf{업데이트:} $F_{RF} \gets [F_{RF}, \mathbf{a}_k]$
            \State \textbf{디지털 빔포머 계산:} $F_{BB} \gets (F_{RF}^H F_{RF})^{-1} F_{RF}^H F_{opt}$
            \State \textbf{잔여 행렬 업데이트:} $R \gets F_{opt} - F_{RF}F_{BB}$
        \EndFor
        \State \textbf{출력:} $F_{RF}, F_{BB}$
    \end{algorithmic}
\end{algorithm}

\subsection{핵심 한계}
채널의 특이값($\sigma_i$)이 큰, 즉 채널 용량에 더 크게 기여하는 중요한 스트림과 그렇지 않은 스트림을 구분하지 않고 모든 스트림의 오차를 동일하게 취급합니다.

\section{SWOMP (Simultaneous Weighted OMP)}
SWOMP는 표준 OMP의 한계를 극복하기 위해 각 스트림의 중요도(채널 특이값)를 고려한 알고리즘입니다. 여기서 \textbf{`Simultaneous(동시)`}라는 키워드는 단 하나의 아날로그 빔포머 $F_{RF}$가 모든 스트림에 \textbf{공유되어 동시에 적용}된다는 하이브리드 빔포밍의 하드웨어적 제약조건을 반영합니다.

\subsection{목표 함수}
가중치 행렬 $W = \text{diag}(\sigma_1, \dots, \sigma_{N_s})$를 도입하여 가중된 근사 오차를 최소화합니다.
\begin{equation}
    \min_{F_{RF}, F_{BB}} ||(F_{opt} - F_{RF}F_{BB})W||_F^2
\end{equation}
이 목표 함수는 $\sigma_i$가 큰 스트림 $i$에서 발생하는 오차를 더 중요하게 다루도록 유도합니다.

\subsection{핵심 차이점: 가중 매칭 (Weighted Matching)}
OMP와의 핵심 차이점은 빔을 선택하는 매칭 단계에 있습니다. 단순히 잔여 행렬 $R$이 아닌 \textbf{가중된 잔여 행렬 ($R \cdot W$)}을 사용합니다. 이는 선택될 빔이 가중치가 높은 중요한 스트림에 미치는 영향을 우선적으로 고려하도록 만듭니다.

\subsection{의사 코드}
\begin{algorithm}[H]
    \caption{Simultaneous Weighted OMP (SWOMP)}
    \begin{algorithmic}[1]
        \State \textbf{입력:} $F_{opt} \in \mathbb{C}^{N_t \times N_s}$, 코드북 $\mathcal{A} \in \mathbb{C}^{N_t \times N_c}$, $N_{RF}$, 가중치 $W = \text{diag}(\sigma_1, \dots, \sigma_{N_s})$
        \State \textbf{초기화:} 잔여 행렬 $R \gets F_{opt}$, 아날로그 빔포머 $F_{RF} \gets [\ ]$
        \For{$k = 1$ to $N_{RF}$}
            \State \textbf{\textcolor{red}{가중 매칭(Weighted Matching):}} \textbf{가중된} 잔여 행렬 $R \cdot W$ 와의 상관관계가 가장 큰 빔을 선택
            \begin{equation*}
                \mathbf{a}_k \gets \arg\max_{\mathbf{a} \in \mathcal{A}} ||\mathbf{a}^H (R \cdot W)||_F^2
            \end{equation*}
            \State \textbf{업데이트:} $F_{RF} \gets [F_{RF}, \mathbf{a}_k]$
            \State \textbf{디지털 빔포머 계산:} $F_{BB} \gets (F_{RF}^H F_{RF})^{-1} F_{RF}^H F_{opt}$
            \State \textbf{잔여 행렬 업데이트:} $R \gets F_{opt} - F_{RF}F_{BB}$
        \EndFor
        \State \textbf{출력:} $F_{RF}, F_{BB}$
    \end{algorithmic}
\end{algorithm}

\section{알고리즘 비교 요약}

\begin{table}[H]
\centering
\begin{tabular}{|l|l|l|}
\hline
\textbf{구분} & \textbf{Standard OMP} & \textbf{SWOMP} \\ \hline
\textbf{핵심 아이디어} & 모든 스트림을 동등하게 고려 & 중요 스트림에 가중치 부여 \\ \hline
\textbf{목표 함수} & $||F_{opt} - F_{RF}F_{BB}||_F^2$ & $||(F_{opt} - F_{RF}F_{BB})W||_F^2$ \\ \hline
\textbf{매칭 기준} & $\max ||\mathbf{a}^H R||_F^2$ & $\max ||\mathbf{a}^H (R \cdot W)||_F^2$ \\ \hline
\textbf{핵심 차별점} & 기준 알고리즘 & \textbf{가중된 잔여 행렬}을 매칭에 사용 \\ \hline
\end{tabular}
\caption{OMP와 SWOMP 비교}
\end{table}

결론적으로 SWOMP는 OMP의 한계를 극복하기 위해 채널 특이값을 가중치로 활용하여, 제한된 RF 자원을 총 처리율(Sum-rate)에 가장 큰 영향을 미치는 주요 고유 모드에 집중시켜 시스템 전체의 성능을 극대화하는 알고리즘입니다.

\end{document}